% ============================================================
% Chapter 4 — Methodology
% ============================================================

\chapter{Methodology}
\label{chap:methodology}

This chapter describes the experimental ladder---a progression of models that isolates the contribution of graph structure, heterogeneous typing, and self-supervised pretraining to fraud detection.

% ────────────────────────────────────────────────────────────
\section{Experimental Ladder}
\label{sec:experimental-ladder}

\begin{description}[style=nextline]
  \item[L0 — Tabular Baselines]
    Logistic Regression and XGBoost on transaction features only.
  \item[L1 — Graph-Augmented Features]
    Hand-crafted graph features fed to XGBoost.
  \item[L2 — Homogeneous GNNs and KGE]
    GCN, GraphSAGE, TransE, and DistMult on a homogeneous graph.
  \item[L3 — Heterogeneous GNNs]
    HGT and HMPNN on typed graph variants (V1, V2, TXN\_V1).
  \item[L4 — Self-Supervised Pretraining]
    HGMAE and LaundroGraph pretraining, then supervised fine-tuning.
\end{description}

Each transition answers a specific question:
L0$\to$L1: does graph structure help?
L1$\to$L2: do learned representations outperform hand-crafted features?
L2$\to$L3: does heterogeneous typing improve over homogeneous?
L3$\to$L4: does self-supervised pretraining add value under label scarcity?

% ────────────────────────────────────────────────────────────
\section{L0: Tabular Baselines}
\label{sec:l0}

\paragraph{Logistic Regression.}
An $\ell_2$-regularised linear classifier with inverse-frequency class weighting, operating on the edge feature vector described in \Cref{sec:feature-engineering}.

\paragraph{XGBoost.}
A gradient-boosted tree ensemble on the same features, with \texttt{scale\_pos\_weight} set to the inverse class frequency ratio.

% ────────────────────────────────────────────────────────────
\section{L1: Graph-Augmented Features}
\label{sec:l1}

The tabular feature vector is augmented with graph statistics computed from the training-period graph: sender/receiver degree, per-account amount statistics (mean, std, total), counterparty diversity, channel entropy, and temporal behaviour ratios (night/weekend proportions). These are concatenated with the edge features and fed to XGBoost.

% ────────────────────────────────────────────────────────────
\section{L2: Homogeneous Models}
\label{sec:l2}

\paragraph{GCN.}
Spectral convolution \citep{kipf2017gcn} with layer-wise propagation:
\begin{equation}
  \mathbf{H}^{(l+1)} = \sigma\!\left(\tilde{\mathbf{D}}^{-\frac{1}{2}} \tilde{\mathbf{A}} \tilde{\mathbf{D}}^{-\frac{1}{2}} \mathbf{H}^{(l)} \mathbf{W}^{(l)}\right)
\end{equation}
where $\tilde{\mathbf{A}} = \mathbf{A} + \mathbf{I}$. For edge classification, source and target embeddings are concatenated with the edge feature vector and passed through an MLP.

\paragraph{GraphSAGE.}
Inductive neighbourhood aggregation \citep{hamilton2017sage}:
\begin{equation}
  \mathbf{h}_v^{(l+1)} = \sigma\!\left(\mathbf{W}^{(l)} \cdot \mathrm{CONCAT}\!\left(\mathbf{h}_v^{(l)},\; \mathrm{AGG}\!\left(\left\{\mathbf{h}_u^{(l)} : u \in \mathcal{N}(v)\right\}\right)\right)\right)
\end{equation}
with a mean aggregator. Its inductive design suits the temporal split, where test nodes may be unseen during training.

\paragraph{TransE.}
Translational embedding \citep{bordes2013transe}: $\mathbf{h} + \mathbf{r} \approx \mathbf{t}$ for valid triples, trained with margin-based ranking loss.

\paragraph{DistMult.}
Bilinear scoring \citep{yang2015distmult}: $f(h,r,t) = \mathbf{h}^\top \mathrm{diag}(\mathbf{r})\,\mathbf{t}$, trained with a contrastive objective.

% ────────────────────────────────────────────────────────────
\section{L3: Heterogeneous GNNs}
\label{sec:l3}

\subsection{Graph Topology Variants}
\label{sec:graph-variants}

Three heterogeneous graph constructions are evaluated, each encoding progressively finer type information:

\paragraph{V1 — Onus vs.\ External (2 edge types).}
\begin{align*}
  \texttt{InternalAccount} &\xrightarrow{\texttt{onus\_transfer}} \texttt{InternalAccount} \\
  \texttt{InternalAccount} &\xrightarrow{\texttt{external\_transfer}} \texttt{ExternalAccount}
\end{align*}
Transactions are edges; the single distinction is whether sender and receiver share the same bank.

\paragraph{V2 — Payment-Rail Typed (6 edge types).}
Refines external transfers by payment rail:
\begin{align*}
  \texttt{InternalAccount} &\xrightarrow{\texttt{onus\_transfer}} \texttt{InternalAccount} \\
  \texttt{InternalAccount} &\xrightarrow{\texttt{ext\_realtime}} \texttt{ExternalAccount} \\
  \texttt{InternalAccount} &\xrightarrow{\texttt{ext\_giro}} \texttt{ExternalAccount} \\
  \texttt{InternalAccount} &\xrightarrow{\texttt{ext\_future}} \texttt{ExternalAccount} \\
  \texttt{InternalAccount} &\xrightarrow{\texttt{ext\_salary}} \texttt{ExternalAccount} \\
  \texttt{InternalAccount} &\xrightarrow{\texttt{ext\_other}} \texttt{ExternalAccount}
\end{align*}
This tests whether payment-rail-specific message functions improve detection.

\paragraph{TXN\_V1 — Transactions as Nodes (node classification).}
\begin{align*}
  \texttt{InternalAccount} &\xrightarrow{\texttt{sends}} \texttt{Transaction} \\
  \texttt{Transaction} &\xrightarrow{\texttt{received\_by\_internal}} \texttt{InternalAccount} \\
  \texttt{Transaction} &\xrightarrow{\texttt{received\_by\_external}} \texttt{ExternalAccount}
\end{align*}
Transactions become nodes with their own feature vectors, enabling direct node classification and richer message passing between accounts and transactions.

\subsection{Heterogeneous Graph Transformer (HGT)}

HGT \citep{hu2020hgt} uses type-dependent projections for attention:
\begin{equation}
  \mathbf{Q}_i = \mathbf{W}^Q_{\tau(t)} \mathbf{h}_t, \quad
  \mathbf{K}_i = \mathbf{W}^K_{\tau(s)} \mathbf{h}_s, \quad
  \mathbf{V}_i = \mathbf{W}^V_{\tau(s)} \mathbf{h}_s
\end{equation}
where $\tau(\cdot)$ denotes node type. Attention is computed per meta-relation $(\tau(s), \phi(e), \tau(t))$, learning type-specific importance. For edge classification, source and target embeddings are concatenated with edge features and classified by an MLP.

\subsection{HMPNN}

The Heterogeneous Message Passing Neural Network \citep{johannessen2023hmpnn} uses edge-type-specific message functions incorporating edge features:
\begin{equation}
  \mathbf{m}_{s \to t} = \mathrm{MLP}_{\phi(e)}\!\left(\mathrm{CONCAT}\!\left(\mathbf{h}_s,\; \mathbf{e}_{s,t}\right)\right)
\end{equation}
Messages are mean-pooled per edge type, then summed across types to update target embeddings.

% ────────────────────────────────────────────────────────────
\section{L4: Self-Supervised Pretraining}
\label{sec:l4}

\subsection{Motivation}
\label{sec:l4-motivation}

Confirmed fraud accounts for approximately \todo{FILL}\% of transactions, so the supervised loss is dominated by the legitimate class and the encoder receives meaningful fraud gradient on only a small fraction of steps. Self-supervised pretraining lets the encoder learn from \emph{all} transactions equally before fraud-specific fine-tuning.

Three lines of prior work motivate this level.
\citet{rao2022xfraud} validate HGT as an encoder for production-scale fraud detection.
Autoencoder-based fraud detection \citep{pumsirirat2018autoencoder, singh2024hetero} shows that reconstruction objectives capture anomalous patterns---fraudulent transactions produce higher reconstruction error.
\citet{tian2023hgmae} combine masked autoencoding with heterogeneous structure, achieving strong node classification via type-aware masking.
L4 combines these ideas: HGMAE-style pretraining on a heterogeneous transaction graph, then supervised fine-tuning. The ladder design ensures that any gain over L3 with the \emph{same} encoder is attributable to pretraining alone.

\subsection{HGMAE}

The Heterogeneous Graph Masked Autoencoder \citep{tian2023hgmae} optimises three objectives: (1) metapath-based edge reconstruction with semantic attention, (2) masked node feature restoration with adaptive masking, and (3) positional feature prediction via metapath2vec encodings. After pretraining, encoder weights initialise a downstream classifier fine-tuned on labelled fraud data.

\subsection{LaundroGraph}

Following \citet{cardoso2022laundrograph}, a GAT encoder learns embeddings via self-supervised link prediction on held-out edges. The resulting representations serve as input to a downstream classifier.

\todo{FILL: finalise L4 implementation details as they stabilise.}

% ────────────────────────────────────────────────────────────
\section{Training Protocol}
\label{sec:training-protocol}

All neural models share a common protocol:
\begin{itemize}[nosep]
  \item Adam optimiser ($\beta_1 = 0.9$, $\beta_2 = 0.999$), learning rate \todo{FILL}
  \item Binary cross-entropy with inverse-frequency class weighting
  \item Early stopping: patience \todo{FILL} epochs on validation PR-AUC
  \item Full-batch training (graph fits in memory)
\end{itemize}

% ────────────────────────────────────────────────────────────
\section{Evaluation Framework}
\label{sec:eval-framework}

The primary metric is \textbf{PR-AUC}, which directly measures the precision--recall tradeoff under extreme class imbalance. Secondary metrics: AUC-ROC, F1 at optimal threshold, and false positives per true positive (operational cost). All metrics are computed on the temporally held-out test split.
